\clearpage
\hypertarget{native-473-lower}{}
\begin{center}
\rule{0.55\linewidth}{0.4pt}\\[2pt]
Astronomische Nachrichten. Vol. XXII. No. 516. pp. 189..192. 3 October 1844.\\[2pt]
\rule{0.55\linewidth}{0.4pt}
\end{center}
\begin{flushright}1 September 1844.\end{flushright}

Last night, after a long interruption, I was able to make a couple more comparisons of the comet with neighbouring fixed stars. They could only be used to determine the right ascension

\begin{center}
1844 August 31.\quad $8^{\rm u}\ 25^{\rm m}\ 57^{\rm s}$ mean time\quad R.A. $203^\circ\ 45'\ 8^{\rm u}3$
\end{center}

Dr. \textsc{Goldschmidt}’s parabola still agrees with it to within one second.

I should also mention that the radius vector still has to be included as a factor in \textsc{Goldschmidt}’s formulas for calculating the coordinates as printed in no. 512 of A. N.

\begin{flushright}Göttingen, Sept. 3.\end{flushright}

I am sending you another comet observation.

\begin{center}
1844 Sept. 2.\quad $8^{\rm u}\ 27^{\rm m}\ 11^{\rm s}$\quad R.A. $=203^\circ\ 11'\ 0''2$\\[2pt]
Decl. $=+7^\circ\ 6'\ 32''1$
\end{center}

It is based on two comparisons with a star of the 9th magnitude from \textsc{Bessel}’s zones 83 and 160 and with the two \textsc{Piazzi} stars XIII, 194 and 208. --- The right ascension of the second star (XIII, 194) gave me some cause for doubt, and it almost seems that the star has a proper motion. The time difference between the second and third stars is given as follows

\begin{center}
\setlength{\tabcolsep}{3pt}
\begin{tabular}{lll}
in the Hist. Céleste & $2^{\rm m}\ 26^{\rm s}00$ & \\
in \textsc{Zach}’s catalogue & $2\ 26.11$ & \\
in \textsc{Piazzi}’s cat., 1st ed. & $2\ 26.04$ & each star 4 times \\
\hphantom{in \textsc{Piazzi}’s cat., }2nd ed. & $2\ 26.31$ & each star observed 10 times \\
in \textsc{Bessel}’s zone 160. & $2\ 27.00$ & \\
my two transits give & $2\ 27.64$ & \\
\end{tabular}
\end{center}

\Needspace{6\baselineskip}
\SourceLeaf{474}
The difference therefore seems to be increasing, even though the annual precession for XIII. 208 is a couple of hundredths of an arcsecond smaller than for XIII, 194. --- To determine the comet’s position above, I used only \textsc{Bessel}’s observations of the three stars and took the mean of all the results. If I excluded the result from XIII, 194 when determining the right ascensions, I would obtain
\[
203^\circ\ 10'\ 57''2
\]
\noindent Refraction, incidentally, was carefully taken into account in the reduction. --- \textsc{Goldschmidt}’s parabola still gives a good right ascension, but the declination it gives is 1 minute too small.

\begin{flushright}Göttingen, Sept. 6.\end{flushright}

Yesterday I observed the comet again, but I was able to observe only one transit, which could be used only for the \emph{declination}. The result observed with the right ascension taken into account was

\begin{center}
1844 Sept. 5.\quad $7^{\rm u}\ 56^{\rm m}\ 56^{\rm s}$ mean time\quad Decl. of comet $+5^\circ\ 21'\ 20''2$
\end{center}

\textsc{Goldschmidt}’s parabola gives $1'\ 5''$ less.

\begin{flushright}Göttingen, Sept. 8.\end{flushright}

Today I am sending you my comet observation from the day before yesterday:

\begin{center}
\setlength{\tabcolsep}{3pt}
\begin{tabular}{lll}
1844 Sept. 6. & $8^{\rm u}\ 1^{\rm m}\ 53^{\rm s}$ & Right ascension $202^\circ\ 7'\ 6''7$ \\
 & $8\ 8\ 7$ & Northern declination $4\ 46\ 25.6$ \\
\end{tabular}
\end{center}

The right ascension is based on four comparisons and the declination on two.

Last night the sky was overcast.
\begin{flushright}C. F. \textsc{Gauss}.\end{flushright}
